% !PRINTED 174
\weberpage{174}
Wir betrachten nun als Rationalitätsbereich $\Omega$ den Körper,
der aus den absoluten Invarianten der Gruppe $S$ und allen
Zahlen\footnote{Man kann sich auch auf einen besonderen Zahlkörper beschränken, wenn nur darin die Substitutionsco\"efficienten der Gruppe $S$ enthalten sind.} besteht. Diesem Rationalitätsbereich gehören die Co\"efficienten
von $\Phi(t)$ an, und wir beweisen zunächst den Satz:

\gesperrt{3. Die Function $\Phi(t)$ ist in $\Omega$ irreducibel.}

Ist nämlich $\Psi(t)$ irgend eine Function in $\Omega$, die für $t = \Theta$
verschwindet, so können wir in der Gleichung $\Psi(\Theta) = 0$, da
die $x_1, x_2, \dots, x_n$ unabhängige Variable sind, das System $(x)$
dieser Variablen durch $A(x)$ ersetzen, wenn $A$ irgend eine Substitution
aus $S$ ist. Dadurch kann $\Theta$ in jede der Functionen
$\Theta_1, \Theta_2, \dots, \Theta_{\mu-1}$ übergeführt werden, während die Co\"efficienten
von $\Psi$ ungeändert bleiben, und folglich ist $\Psi(\Theta_1) = 0$,
$\Psi(\Theta_2) = 0, \dots \Psi(\Theta_{\mu-1}) = 0$. Es muss also $\Psi(t)$ durch $\Phi(t)$
theilbar sein, wodurch die Irreducibilität erwiesen ist.

Es bedeute nun $\omega$ irgend eine Function der Variablen $x$ und
\[ \omega,\ \omega_1,\ \omega_2,\ \dots,\ \omega_{\mu-1} \]
mögen die Functionen sein, die aus $\omega$ durch Anwendung der
Substitutionen von $S$ entstehen, von denen nun nicht vorausgesetzt
zu werden braucht, dass sie alle von einander verschieden
sind. Ist $t$ eine Variable, so ist
\weq{7}{\Phi(t)\left(\frac{\omega}{t-\Theta} + \frac{\omega_1}{t-\Theta_1} + \cdots + \frac{\omega_{\mu-1}}{t-\Theta_{\mu-1}}\right) = \Psi(t)}
eine ganze rationale Function $(\mu-1)^{\text{ten}}$ Grades von $t$, und zugleich
ist es eine Invariante von $S$, also eine Function in $\Omega$.
Setzen wir darin $t = \Theta$, so folgt durch einen schon früher oft
angewandten Schluss (Bd.\ I, \S.\,143, 155):
\weq{8}{\omega = \frac{\Psi(\Theta)}{\Phi'(\Theta)},}
worin der Satz enthalten ist:

\gesperrt{4. Jede rationale Function der Variablen $(x)$ kann
rational durch $\Theta$ ausgedrückt werden, gehört
also dem Körper $\Omega(\Theta)$ an.}

Aus diesem Satze können wir einen zweiten Beweis dafür
ableiten, dass alle Invarianten der Gruppe $S$ rational durch eine
endliche Anzahl von ihnen, nämlich die Co\"efficienten der Function
$\Phi(t)$, darstellbar sind. Denn wenn wir eine absolute In-
%CONT>>
